% !TeX root = ../../Book3.tex
\section{二连通区域}
\label{b3:sec:0801}

单连通真平面区域只有圆盘这一内部共形模型。二连通区域开始出现一个不能靠归一化消去的实参数。先区分标准模型与退化类型，再以曲线族给这个参数一个可计算的定义。

\subsection{环域}
\label{b3:subsec:080101}

\begin{definition}\label{b3:def:08-doubly-connected}
扩充平面中区域 \(U\) 称为二连通，若 \(\widehat{\mathbb C\kern0pt}\setminus U\) 恰有两个非空连通分支。每个分支是一个紧连通集，允许退化为单点。
\end{definition}
本章“环域”在广义讨论中指这种区域，标准圆环则另加“圆”字。二连通并不要求两条光滑边界；裂缝、线段或复杂紧连通集都可作为补集分支。拓扑上的环与共形标准圆环也不同：拓扑等价不控制两端之间的解析距离。

\ProseParagraphBreak
\(\mathbb C\setminus\{0\}\) 在球面中删去了0和无穷远，是二连通；单位去心盘的两个补集分支是0与闭外部。普通环 \(r<|z|<R\) 的两个补集分支都是非退化的。要在球面中计算补集，才能给无界区域作一致分类。

\subsection{标准环域}
\label{b3:subsec:080102}

标准圆环记为 \(A(r,R)=\{r<|z|<R\}\)，其中 \(0<r<R<\infty\)。缩放可令外半径为1。映射 \(z\mapsto rR/z\) 交换两端，旋转保持每一端。去心盘 \(\mathbb D^*\) 与去心平面 \(\mathbb C^*\) 是两个退化模型。

\ProseParagraphBreak
取对数后，圆环的万有覆盖是竖直带 \(\log r<\operatorname{Re}w<\log R\)，甲板平移为 \(w\mapsto w+2\pi\mathrm{i}\)。这是直接由指数的周期验证的模型，不需要先使用一般万有覆盖定理。带宽与平移长度的比将成为模量。去心盘对应半平面 \(\operatorname{Re}w<0\)，去心平面对应整个 \(w\) 平面。

\subsection{分类定理}
\label{b3:subsec:080103}

\begin{theorem}[深层二连通分类]\label{b3:thm:08-annulus-classification}
二连通区域 \(U\) 的两个补集分支若都非单点，则 \(U\) 双全纯等价于唯一的 \(A(r,1)\)、\(0<r<1\)。若恰有一个分支为点，则等价于 \(\mathbb D^*\)；若两个都为点，则等价于 \(\mathbb C^*\)。
\end{theorem}
\begin{proofsketch}
任意边界下的存在部分是本章明确采用的深层输入。可先建立万有覆盖与统一化，将二连通域写成单连通模型对无限循环甲板群的商；圆盘模型中的生成元有抛物或双曲类型，其商分别是去心盘与有限圆环，平面模型的平移商为去心平面。还要证明点状补集分支与退化端一致，不能仅凭基本群为整数群区分三者。此路线在第11--13章解释；一般二连通映射可参阅 Ahlfors《Conformal Invariants》\cite{Ahlfors2010Conformal}关于环域与极值长度的论述。标准参数的唯一性在下文由完整的极值长度计算证明，未作为未经说明的假定。
\end{proofsketch}
分类的存在性并不意味着给出显式映射。对于两条不同心圆之间的区域，分式变换可先把两圆同心化；对于一般裂缝环域，常需使用数值或势论方法。

\subsection{共形不变量}
\label{b3:subsec:080104}

有限圆环的半径比是一个真正的共形不变量。同样面积或同样欧氏厚度的两个圆环可有不同半径比，因此不共形等价。另一方面，\(A(r,R)\) 与 \(A(cr,cR)\) 对任意 \(c>0\) 等价，绝对大小并非不变量。

\ProseParagraphBreak
退化类型还需单独记录。去心盘与去心平面都具有下文意义的无限模量，却不等价：若有 \(\mathbb C^*\rightarrow\mathbb D^*\) 双全纯映射，作为有界函数它在0可去，成为有界整函数，必常值。无限模量不能独自区分这两个端部结构。

\ProseParagraphBreak
“二连通”在这里以球面补集的连通分支数定义。对于 \(A(r,R)\)，两个补分支分别是内闭盘和外闭盘加无穷远，都非单点；去心盘的内分支退化为一个点，外分支仍非单点；去心平面的两分支都为点。圆盘本身在球面中只有一个补分支，虽在某幅平面图里可以同时看到“外面”和“无穷远”，它们属于同一个分支，不能重复计数。

\ProseParagraphBreak
两端的拓扑资料还可从穷竭看出。去心盘中离穿孔越来越近的一列小圆与离外圆越来越近的一列大圆代表两个不同端；前者在球面闭包收缩到一个点，后者的极限集是整个外边界。双全纯映射保持内部端的趋近结构，却不必把任意粗糙普通边界逐点对应。这是正文用端部曲线族定义模量、而没有在一般分类里提前假定边界同胚的原因。
